% % % % % % % % % % % % % % % % % % % % % % % % % % % % % % % % % % % % % % % % % % % % % % % % % % % % % % % % % % % % % % % %
% SAT_STYLE_MATH_MCQS.tex
% 1_02.tex
% 
% encoding: UTF-8 
% EOL: LF
%
% format: LaTeX
% indent: spaces (2)
% width: 127
% % % % % % % % % % % % % % % % % % % % % % % % % % % % % % % % % % % % % % % % % % % % % % % % % % % % % % % % % % % % % % % %
Le chiffre des unités pour l'écriture décimale  $3^{100}$ est:
%
\begin{enumerate}
  \item $3$
  \item $9$
  \item $1$
  \item $7$
  \item $6$
\end{enumerate}
%
\begin{proof} %[\textbf{C}]
Rappelons que la division euclienne d'un entier positif $N$ par $10$ retourne $\mathtt{N\,\% \, 10}$, l'entier naturel $< 10$ % 
qui satisfait 
%
\begin{equation}
  \mathtt{N = (N \, //\, 10 ) \ast 10  +(N \, \% \, 10)}.
\end{equation}
%
Rechercher le chiffre des unités revient donc à trouver le \textit{reste} $\mathtt{N\,\% \, 10}$. Pour ce faire, %
partons de 
%
\begin{equation}
  3^{100} = (3^2)^{50} = 9^{50} =(9^2)^{25}
\end{equation}
%
puis remarquons que 
\begin{equation}
  9^2 = 81 \equiv 1 \mod 10.
\end{equation}
%
D'où
%
\begin{equation}
  3^{100} = (9^2)^{25} \equiv 1^{25} \equiv 1 \mod 10.
\end{equation}
%
Alternativement, on peut aussi noter que 
%
\begin{equation}
  3^{100} = 9^{50} \equiv (\minus 1)^{50} \equiv 1 \mod 10. 
\end{equation}
%
Réponse \textbf{C}.
%
\end{proof}